\chapter{Theory of Generalization}
\lecture{6}{30 Sep.}{}

\begin{prev}
    \autoref{sec:break-point} ended with the growth function
    \(m_{\mathcal{H}}(N)\), the number of dichotomies \(\mathcal{H}\) can
    implement on the best \(N\) inputs, as the (wishful) effective price of
    choice in training, and with \autoref{conj:break}: a break point \(k\)
    should force \(m_{\mathcal{H}}(N) = O(N^{k-1})\). Two things were left
    open --- proving that conjecture, and removing the `\(?\)' from
    \[
        \mathbb{P}\left[ \abs{\Ein(g) - \Eout(g)} > \epsilon \right]
        \overset{?}{\le} 2 \cdot m_{\mathcal{H}}(N) \cdot
        \exp \left( -2 \epsilon^2 N \right) .
    \]
\end{prev}

This chapter does both. The first three sections prove the conjecture without
ever looking at a particular \(\mathcal{H}\): a break point is a purely
combinatorial restriction on which tables of \(\circ\) and \(\times\) may occur,
and counting the largest such table gives a polynomial. The last section
sketches why \(m_{\mathcal{H}}\), with a few changes of constant, really can
stand in for \(M\).

\section{Restriction of Break Point}\label{sec:restriction}

\source{slides p.2--5}

\subsection{The Four Break Points}

The four hypothesis sets of the last chapter, each with the smallest input
configuration it fails to shatter:

\begin{center}
    \solidrules
    \begin{adjustbox}{max width=\textwidth}
        \begin{tabular}{@{}l|c|c|l@{}}
            \toprule
            \(\mathcal{H}\) & \(m_{\mathcal{H}}(N)\) & inputs not shattered
            & break point\\
            \midrule
            positive rays & \(N + 1\) & \dmark{o}\,\dmark{x}
            & \(m_{\mathcal{H}}(2) = 3 < 2^2\): break point at \(2\)\\
            positive intervals & \(\tfrac{1}{2} N^2 + \tfrac{1}{2} N + 1\)
            & \dmark{o}\,\dmark{x}\,\dmark{o}
            & \(m_{\mathcal{H}}(3) = 7 < 2^3\): break point at \(3\)\\
            convex sets & \(2^N\)
            & \tikz[baseline=(current bounding box.center), x=0.3cm, y=0.3cm]{%
                \dmarkat{1}{2}{o}\dmarkat{2}{2}{x}\dmarkat{0}{1}{o}%
                \dmarkat{1}{0}{x}\dmarkat{2}{0}{o}}
            & \(m_{\mathcal{H}}(N) = 2^N\) always: no break point\\
            2D perceptrons & \(< 2^N\) in some cases
            & \tikz[baseline=(current bounding box.center), x=0.3cm, y=0.3cm]{%
                \dmarkat{1}{2}{o}\dmarkat{0}{1}{x}\dmarkat{2}{1}{x}%
                \dmarkat{1}{0}{o}}
            & \(m_{\mathcal{H}}(4) = 14 < 2^4\): break point at \(4\)\\
            \bottomrule
        \end{tabular}
    \end{adjustbox}
\end{center}

The third column is the dichotomy that fails: a ray cannot say \(\circ\) left
of \(\times\), an interval cannot leave a hole in the middle, a line cannot do
XOR. Convex sets never fail, whatever the input.

\autoref{sec:break-point} already showed one consequence of a break point:
break point \(k\) \(\implies\) break point \(k + 1, \dots\) That concerns only
the values \(m_{\mathcal{H}}(N) < 2^N\) at \(N \ge k\).

\begin{question}\label{q:what-else}
    What else does a break point imply?
\end{question}

\subsection{What Must Be True When \texorpdfstring{\(k = 2\)}{k = 2}}

Suppose the \emph{minimum} break point is \(k = 2\). Without knowing
\(\mathcal{H}\), a few values of \(m_{\mathcal{H}}\) are forced:

\begin{itemize}
    \item \(N = 1\): every \(m_{\mathcal{H}}(N) = 2\) by definition, since the
        minimum break point is above \(1\), so some single input is shattered;
    \item \(N = 2\): every \(m_{\mathcal{H}}(N) < 4\) by definition, since \(2\)
        is a break point (so the maximum possible is \(3\)).
\end{itemize}

\begin{question}\label{q:max-n3-k2}
    What is the maximum possible \(m_{\mathcal{H}}(N)\) when \(N = 3\) and
    \(k = 2\)?
\end{question}

The only rule is that no two of the three inputs may be shattered: in the table
of dichotomies, no pair of columns may show all four patterns
\(\circ\circ, \circ\times, \times\circ, \times\times\). So build the table row
by row and reject any row that breaks the rule.

\begin{itemize}
    \item \(\circ\circ\circ\): \(1\) dichotomy; shatters any two points? No.
    \item add \(\circ\circ\times\): \(2\) dichotomies; no.
    \item add \(\circ\times\circ\): \(3\) dichotomies; no.
    \item add \(\circ\times\times\): \(4\) dichotomies, but now columns
        \(\bm{x}_2, \bm{x}_3\) read \(\circ\circ, \circ\times, \times\circ,
        \times\times\) --- shattered. Rejected.
    \item add \(\times\circ\circ\) instead: \(4\) dichotomies; no pair
        shattered.
    \item a fifth row: each of the three remaining candidates shatters some
        pair (table below).
\end{itemize}

\begin{center}
    \solidrules
    \begin{tabular}{@{}ccc|l@{}}
        \toprule
        \(\bm{x}_1\) & \(\bm{x}_2\) & \(\bm{x}_3\) & \\
        \midrule
        \dmark{o} & \dmark{o} & \dmark{o} & kept\\
        \dmark{o} & \dmark{o} & \dmark{x} & kept\\
        \dmark{o} & \dmark{x} & \dmark{o} & kept\\
        \dmark{x} & \dmark{o} & \dmark{o} & kept\\
        \midrule
        \dmark{o} & \cellcolor{shatcell}\dmark{x} & \cellcolor{shatcell}\dmark{x}
        & 4th row tried: shatters \(\bm{x}_2, \bm{x}_3\)\\
        \cellcolor{shatcell}\dmark{x} & \dmark{o} & \cellcolor{shatcell}\dmark{x}
        & 5th row tried: shatters \(\bm{x}_1, \bm{x}_3\)\\
        \cellcolor{shatcell}\dmark{x} & \cellcolor{shatcell}\dmark{x} & \dmark{o}
        & 5th row tried: shatters \(\bm{x}_1, \bm{x}_2\)\\
        \cellcolor{shatcell}\dmark{x} & \cellcolor{shatcell}\dmark{x} & \dmark{x}
        & 5th row tried: shatters \(\bm{x}_1, \bm{x}_2\) (every pair, in fact)\\
        \bottomrule
    \end{tabular}
\end{center}

The four kept rows use up the table: \(\circ\times\times\) was already refused,
and \(\times\circ\times\), \(\times\times\circ\), \(\times\times\times\) each
complete the four patterns on some pair, so no column pair has room left. The
maximum possible so far is \(4\) dichotomies. The search here started from
\(\circ\circ\circ\) and so is not a proof that no other table of five exists;
that proof comes for free from \autoref{thm:bounding} below, which gives
\(B(3, 2) \le 4\).

\subsection{Restriction of Break Point}

Putting the three cases together, when the minimum break point is \(k = 2\):

\begin{itemize}
    \item \(N = 1\): every \(m_{\mathcal{H}}(N) = 2\) by definition;
    \item \(N = 2\): every \(m_{\mathcal{H}}(N) < 4\) by definition (so maximum
        possible \(= 3\));
    \item \(N = 3\): maximum possible \(= 4 \ll 2^3\).
\end{itemize}

At \(N = 3\) the break point already bites twice as hard as the definition
asks: the definition only forbids shattering two points, yet it caps three
points at half of \(2^3\).

\begin{moral}
    A break point \(k\) restricts the maximum possible \(m_{\mathcal{H}}(N)\) a
    lot for \(N > k\).
\end{moral}

This is the answer to \autoref{q:what-else}, and it suggests how to prove
\autoref{conj:break} without knowing anything else about \(\mathcal{H}\):
\begin{equation}\label{eq:idea}
    m_{\mathcal{H}}(N)
    \;\le\; \text{maximum possible } m_{\mathcal{H}}(N) \text{ given } k
    \;\le\; \mathrm{poly}(N) .
\end{equation}
The first inequality is trivial; the work is in the second.

\begin{exercise}[Fun Time]
    When the minimum break point is \(k = 1\), what is the maximum possible
    \(m_{\mathcal{H}}(N)\) when \(N = 3\)?
    \begin{enumerate*}[(1)]
        \item \(1\); \item \(2\); \item \(4\); \item \(8\).
    \end{enumerate*}
\end{exercise}

\begin{proof}[Reference Answer]
    (1). Because \(k = 1\), the hypothesis set cannot even shatter one point.
    Thus, every `column' of the table cannot contain both \(\circ\) and
    \(\times\). Then, after including the first dichotomy, it is not possible
    to include any other different dichotomy:

    \begin{center}
        \solidrules
        \begin{tabular}{@{}ccc|l@{}}
            \toprule
            \(\bm{x}_1\) & \(\bm{x}_2\) & \(\bm{x}_3\) & \\
            \midrule
            \dmark{o} & \dmark{x} & \cellcolor{shatcell}\dmark{o} & kept\\
            \dmark{o} & \dmark{x} & \cellcolor{shatcell}\dmark{x}
            & rejected: column \(\bm{x}_3\) now shattered\\
            \bottomrule
        \end{tabular}
    \end{center}

    Any second row differs from the first somewhere, and that column then shows
    both symbols. Thus, the maximum possible \(m_{\mathcal{H}}(N)\) is \(1\).
\end{proof}

\section{Bounding Function: Basic Cases}\label{sec:bound-basic}

\source{slides p.6--11}

The middle term of \autoref{eq:idea} deserves a name.

\subsection{Bounding Function}

\begin{definition}[Bounding function]\label{def:bounding}
    The \vocab{bounding function} \(B(N, k)\) is the maximum possible
    \(m_{\mathcal{H}}(N)\) when \(k\) is a break point.
\end{definition}

Two remarks make this concrete.

\begin{itemize}
    \item It is a \emph{combinatorial quantity}: the maximum number of
        length-\(N\) vectors with entries in \(\left\{ \circ, \times
        \right\}\) such that no length-\(k\) subvector is shattered --- that
        is, no choice of \(k\) coordinates shows all \(2^k\) patterns.
    \item It is \emph{irrelevant of the details of \(\mathcal{H}\)}. For
        example, \(B(N, 3)\) bounds both positive intervals (\(k = 3\)) and
        1D perceptrons (\(k = 3\)).
\end{itemize}

The first remark is why the second holds: the dichotomies of any
\(\mathcal{H}\) with break point \(k\), on any \(N\) inputs, form such a set
of vectors, so
\[
    m_{\mathcal{H}}(N) \le B(N, k)
\]
for every \(\mathcal{H}\) with break point \(k\). The price is looseness ---
\(B\) takes the worst table, whether or not any \(\mathcal{H}\) realises it.

\begin{question}\label{q:B-poly}
    New goal: is \(B(N, k) \le \mathrm{poly}(N)\)?
\end{question}

\subsection{Table of Bounding Function}

We fill in a table of \(B(N, k)\), rows \(N\), columns \(k\). What is known
so far is:

\begin{itemize}
    \item \(B(2, 2) = 3\) (maximum \(< 4\));
    \item \(B(3, 2) = 4\) (the `pictorial' proof above).
\end{itemize}

Three whole regions of the table then follow from short arguments.

\begin{proposition}[Boundary cases]\label{prop:B-basic}
    \leavevmode
    \begin{enumerate}
        \item \(B(N, 1) = 1\) for every \(N\).
        \item \(B(N, k) = 2^N\) for \(N < k\) --- including all dichotomies does
            not violate the `breaking condition'.
        \item \(B(N, k) = 2^N - 1\) for \(N = k\) --- removing a single dichotomy
            satisfies the `breaking condition'.
    \end{enumerate}
\end{proposition}

\begin{proof}
    (1) is the Fun Time above, for any \(N\): if no single coordinate may show
    both symbols, all rows agree everywhere, so there is one row.

    (2) With \(N < k\) there are no \(k\) coordinates to choose, so the
    condition is empty and all \(2^N\) vectors are allowed.

    (3) With \(N = k\) the only \(k\) coordinates are all of them. All \(2^N\)
    vectors would shatter them, so at most \(2^N - 1\); and deleting any one
    vector leaves the full coordinate set missing a pattern, so \(2^N - 1\) is
    attained.
\end{proof}

Note that (3) contains the first known value, \(B(2, 2) = 2^2 - 1 = 3\). The
four stages of the table, merged into one --- \textcolor{blue!75!black}{blue}:
column \(k = 1\); black: \(N < k\); \textcolor{orange!85!black}{orange}:
\(N = k\); purple: \(B(3, 2)\), the one entry from the pictorial search:

\begin{center}
    \solidrules
    \begin{tabular}{@{}cc|ccccccc@{}}
        \toprule
        & & \multicolumn{7}{c}{\(k\)}\\
        & \(B(N, k)\) & \(1\) & \(2\) & \(3\) & \(4\) & \(5\) & \(6\)
        & \(\dots\)\\
        \midrule
        & \(1\) & \textcolor{blue!75!black}{\(1\)} & \(2\) & \(2\) & \(2\)
        & \(2\) & \(2\) & \(\dots\)\\
        & \(2\) & \textcolor{blue!75!black}{\(1\)}
        & \textcolor{orange!85!black}{\(\bm{3}\)} & \(4\) & \(4\) & \(4\) & \(4\)
        & \(\dots\)\\
        & \(3\) & \textcolor{blue!75!black}{\(1\)} & \textcolor{Purple}{\(4\)}
        & \textcolor{orange!85!black}{\(\bm{7}\)} & \(8\) & \(8\) & \(8\)
        & \(\dots\)\\
        \(N\) & \(4\) & \textcolor{blue!75!black}{\(1\)} & &
        & \textcolor{orange!85!black}{\(\bm{15}\)} & \(16\) & \(16\)
        & \(\dots\)\\
        & \(5\) & \textcolor{blue!75!black}{\(1\)} & & &
        & \textcolor{orange!85!black}{\(\bm{31}\)} & \(32\) & \(\dots\)\\
        & \(6\) & \textcolor{blue!75!black}{\(1\)} & & & &
        & \textcolor{orange!85!black}{\(\bm{63}\)} & \(\dots\)\\
        & \(\vdots\) & \(\vdots\) & & & & & & \(\ddots\)\\
        \bottomrule
    \end{tabular}
\end{center}

(The diagonal's first entry \(B(1, 1) = 2^1 - 1 = 1\) belongs to both the
blue column and the orange diagonal.) Everything on and above the diagonal is
now known, and so is the first column; only the region \(N > k \ge 2\) below
the diagonal remains, which is where \(N\) grows with \(k\) fixed --- the only
part that matters for \autoref{q:B-poly}.

\begin{moral}
    More than halfway done! :-)
\end{moral}

\begin{exercise}[Fun Time]
    For the 2D perceptrons, which of the following claims is true?
    \begin{enumerate}[(1)]
        \item minimum break point \(k = 2\);
        \item \(m_{\mathcal{H}}(4) = 15\);
        \item \(m_{\mathcal{H}}(N) < B(N, k)\) when \(N = k =\) minimum break
            point;
        \item \(m_{\mathcal{H}}(N) > B(N, k)\) when \(N = k =\) minimum break
            point.
    \end{enumerate}
\end{exercise}

\begin{proof}[Reference Answer]
    (3). As discussed previously, the minimum break point for 2D perceptrons
    is \(4\), with \(m_{\mathcal{H}}(4) = 14\) (\autoref{prop:fourteen}).
    Also, note that \(B(4, 4) = 15\). So the bounding function \(B(N, k)\) can
    be `loose' in bounding \(m_{\mathcal{H}}(N)\): lines miss two of the
    sixteen dichotomies of four points (a Radon pair and its mirror), where
    the break-point rule alone forces only one to be missing. (4) is
    impossible, since \(B\) is by definition an upper bound.
\end{proof}

\section{Bounding Function: Inductive Cases}\label{sec:bound-inductive}

\source{slides p.12--20}

\subsection{Estimating \texorpdfstring{\(B(4, 3)\)}{B(4, 3)}}

The first unknown entry below the diagonal is \(B(4, 3)\), marked `\(?\)':

\begin{center}
    \solidrules
    \setlength{\extrarowheight}{0.2ex}
    \begin{tabular}{@{}cc|ccccccc@{}}
        \toprule
        & & \multicolumn{7}{c}{\(k\)}\\
        & \(B(N, k)\) & \(1\) & \(2\) & \(3\) & \(4\) & \(5\) & \(6\)
        & \(\dots\)\\
        \midrule
        & \(1\) & \(1\) & \(2\) & \(2\) & \(2\) & \(2\) & \(2\) & \(\dots\)\\
        & \(2\) & \(1\) & \(3\) & \(4\) & \(4\) & \(4\) & \(4\) & \(\dots\)\\
        & \(3\) & \(1\) & \(4\) & \(7\) & \(8\) & \(8\) & \(8\) & \(\dots\)\\
        \(N\) & \(4\) & \(1\) & & \textcolor{orange!85!black}{\(\bm{?}\)}
        & \(15\) & \(16\) & \(16\) & \(\dots\)\\
        & \(5\) & \(1\) & & & & \(31\) & \(32\) & \(\dots\)\\
        & \(6\) & \(1\) & & & & & \(63\) & \(\dots\)\\
        & \(\vdots\) & \(\vdots\) & & & & & & \(\ddots\)\\
        \bottomrule
    \end{tabular}
\end{center}

The motivation for how to get at it: a table on four inputs is a table on three
inputs with one more column, so \(B(4, 3)\) shall be related to \(B(3, ?)\) ---
`adding' one point to a \(B(3, ?)\) table. The next step is to reduce
\(B(4, 3)\) to \(B(3, ?)\).

\subsection{`Achieving' Dichotomies of \texorpdfstring{\(B(4, 3)\)}{B(4, 3)}}

First the answer, found the brute-force way: after checking all \(2^{2^4}\)
sets of dichotomies on four inputs --- every subset of the \(16\) possible rows
--- the winner is a set of \(11\), and \(B(4, 3) = 11\):

\begin{center}
    \solidrules
    \setlength{\extrarowheight}{0.2ex}
    \begin{tabular}[t]{@{}c|cccc@{}}
        \toprule
        & \(\bm{x}_1\) & \(\bm{x}_2\) & \(\bm{x}_3\) & \(\bm{x}_4\)\\
        \midrule
        01 & \dmark{o} & \dmark{o} & \dmark{o} & \dmark{o}\\
        02 & \dmark{x} & \dmark{o} & \dmark{o} & \dmark{o}\\
        03 & \dmark{o} & \dmark{x} & \dmark{o} & \dmark{o}\\
        04 & \dmark{o} & \dmark{o} & \dmark{x} & \dmark{o}\\
        05 & \dmark{o} & \dmark{o} & \dmark{o} & \dmark{x}\\
        06 & \dmark{x} & \dmark{x} & \dmark{o} & \dmark{x}\\
        \bottomrule
    \end{tabular}
    \quad
    \begin{tabular}[t]{@{}c|cccc@{}}
        \toprule
        & \(\bm{x}_1\) & \(\bm{x}_2\) & \(\bm{x}_3\) & \(\bm{x}_4\)\\
        \midrule
        07 & \dmark{x} & \dmark{o} & \dmark{x} & \dmark{o}\\
        08 & \dmark{x} & \dmark{o} & \dmark{o} & \dmark{x}\\
        09 & \dmark{o} & \dmark{x} & \dmark{x} & \dmark{o}\\
        10 & \dmark{o} & \dmark{x} & \dmark{o} & \dmark{x}\\
        11 & \dmark{o} & \dmark{o} & \dmark{x} & \dmark{x}\\
        \bottomrule
    \end{tabular}
    \qquad
    \begin{tabular}[t]{@{}cc|cccccc@{}}
        \toprule
        & & \multicolumn{6}{c}{\(k\)}\\
        & \(B(N, k)\) & \(1\) & \(2\) & \(3\) & \(4\) & \(5\) & \(6\)\\
        \midrule
        & \(1\) & \(1\) & \(2\) & \(2\) & \(2\) & \(2\) & \(2\)\\
        & \(2\) & \(1\) & \(3\) & \(4\) & \(4\) & \(4\) & \(4\)\\
        \rowcolor{singlerowalt}
        & \(3\) & \(1\) & \(4\) & \(7\) & \(8\) & \(8\) & \(8\)\\
        \(N\) & \(4\) & \(1\) & & \cellcolor{singlerowalt}%
        \textcolor{orange!85!black}{\(\bm{11}\)} & \(15\) & \(16\) & \(16\)\\
        & \(5\) & \(1\) & & & & \(31\) & \(32\)\\
        & \(6\) & \(1\) & & & & & \(63\)\\
        \bottomrule
    \end{tabular}
\end{center}

(I reran the search: no set of \(12\) avoids shattering three coordinates, and
the table above passes.) The shaded cells are the hint for what comes next: the new entry \(11\) is to
be explained by the row \(N = 3\) above it.

\begin{question}\label{q:reduce}
    How to reduce \(B(4, 3)\) to \(B(3, ?)\) cases?
\end{question}

\subsection{Reorganized Dichotomies of \texorpdfstring{\(B(4, 3)\)}{B(4, 3)}}

Sort the \(11\) rows by their first three entries, and look at what
\(\bm{x}_4\) does. Some prefixes \((\bm{x}_1, \bm{x}_2, \bm{x}_3)\) occur twice,
once with \(\bm{x}_4 = \circ\) and once with \(\bm{x}_4 = \times\); the others
occur once.

\begin{center}
    \solidrules
    \setlength{\extrarowheight}{0.2ex}
    \begin{tabular}{@{}c|ccc|c@{}}
        \toprule
        & \(\bm{x}_1\) & \(\bm{x}_2\) & \(\bm{x}_3\) & \(\bm{x}_4\)\\
        \midrule
        \rowcolor{pairrow} 01 & \dmark{o} & \dmark{o} & \dmark{o} & \dmark{o}\\
        \rowcolor{pairrow} 05 & \dmark{o} & \dmark{o} & \dmark{o} & \dmark{x}\\
        \rowcolor{pairrowalt} 02 & \dmark{x} & \dmark{o} & \dmark{o} & \dmark{o}\\
        \rowcolor{pairrowalt} 08 & \dmark{x} & \dmark{o} & \dmark{o} & \dmark{x}\\
        \rowcolor{pairrow} 03 & \dmark{o} & \dmark{x} & \dmark{o} & \dmark{o}\\
        \rowcolor{pairrow} 10 & \dmark{o} & \dmark{x} & \dmark{o} & \dmark{x}\\
        \rowcolor{pairrowalt} 04 & \dmark{o} & \dmark{o} & \dmark{x} & \dmark{o}\\
        \rowcolor{pairrowalt} 11 & \dmark{o} & \dmark{o} & \dmark{x} & \dmark{x}\\
        \rowcolor{singlerow} 06 & \dmark{x} & \dmark{x} & \dmark{o} & \dmark{x}\\
        \rowcolor{singlerowalt} 07 & \dmark{x} & \dmark{o} & \dmark{x} & \dmark{o}\\
        \rowcolor{singlerow} 09 & \dmark{o} & \dmark{x} & \dmark{x} & \dmark{o}\\
        \bottomrule
    \end{tabular}
\end{center}

Orange: pair; purple: single. Write \(\alpha\) for the number of pairs and
\(\beta\) for the number of singles. Here \(\alpha = 4\) and \(\beta = 3\), and
\[
    B(4, 3) = 11 = 2\alpha + \beta .
\]
The point of the split is that the two kinds of prefix face different
restrictions, and each restriction is a smaller bounding function.

\subsection{Estimating Part of \texorpdfstring{\(B(4, 3)\)}{B(4, 3)}}

Delete \(\bm{x}_4\) and merge each pair into one row. What remains is a table
on \((\bm{x}_1, \bm{x}_2, \bm{x}_3)\) with \(\alpha + \beta\) distinct rows:

\begin{center}
    \solidrules
    \begin{tabular}[t]{@{}c|ccc@{}}
        \toprule
        & \(\bm{x}_1\) & \(\bm{x}_2\) & \(\bm{x}_3\)\\
        \midrule
        \rowcolor{pairrow} & \dmark{o} & \dmark{o} & \dmark{o}\\
        \rowcolor{pairrowalt} & \dmark{x} & \dmark{o} & \dmark{o}\\
        \rowcolor{pairrow} & \dmark{o} & \dmark{x} & \dmark{o}\\
        \rowcolor{pairrowalt} \multirow{-4}{*}{\(\alpha\)}
        & \dmark{o} & \dmark{o} & \dmark{x}\\
        \rowcolor{singlerow} & \dmark{x} & \dmark{x} & \dmark{o}\\
        \rowcolor{singlerowalt} & \dmark{x} & \dmark{o} & \dmark{x}\\
        \rowcolor{singlerow} \multirow{-3}{*}{\(\beta\)}
        & \dmark{o} & \dmark{x} & \dmark{x}\\
        \bottomrule
    \end{tabular}
    \qquad\qquad
    \begin{tabular}[t]{@{}c|ccc@{}}
        \toprule
        & \(\bm{x}_1\) & \(\bm{x}_2\) & \(\bm{x}_3\)\\
        \midrule
        \rowcolor{pairrow} & \dmark{o} & \dmark{o} & \dmark{o}\\
        \rowcolor{pairrowalt} & \dmark{x} & \dmark{o} & \dmark{o}\\
        \rowcolor{pairrow} & \dmark{o} & \dmark{x} & \dmark{o}\\
        \rowcolor{pairrowalt} \multirow{-4}{*}{\(\alpha\)}
        & \dmark{o} & \dmark{o} & \dmark{x}\\
        \bottomrule
    \end{tabular}
\end{center}

\emph{Left: all \(\alpha + \beta\) prefixes. Right: the \(\alpha\) paired ones
alone.}

\paragraph{The whole prefix table (1/2).}

\begin{itemize}
    \item \(\alpha + \beta\): dichotomies on \((\bm{x}_1, \bm{x}_2, \bm{x}_3)\).
    \item \(B(4, 3)\) `no shatter' any \(3\) inputs \(\implies\) \(\alpha +
        \beta\) `no shatter' any \(3\).
\end{itemize}

The implication holds because deleting a column cannot create a pattern: if
the prefix table showed all \(8\) patterns on \(\bm{x}_1, \bm{x}_2, \bm{x}_3\),
so would the original \(11\) rows. Hence

\begin{moral}
    \(\alpha + \beta \le B(3, 3)\).
\end{moral}

\paragraph{The paired prefixes (2/2).}

\begin{itemize}
    \item \(\alpha\): dichotomies on \((\bm{x}_1, \bm{x}_2, \bm{x}_3)\) with
        \(\bm{x}_4\) \emph{paired}.
    \item \(B(4, 3)\) `no shatter' any \(3\) inputs \(\implies\) \(\alpha\)
        `no shatter' any \(2\).
\end{itemize}

The claim is about the \(\alpha\) rows alone, not the whole table: rows 01,
02, 03, 06 already shatter \(\bm{x}_1, \bm{x}_2\), and a break point at \(3\)
says nothing about \(2\) in general. What makes the \(\alpha\) rows special is
that each comes with \emph{both} values of \(\bm{x}_4\). So if they shattered
two inputs, say \(\bm{x}_1, \bm{x}_2\), the original table would show each of
those four patterns with \(\bm{x}_4 = \circ\) and with \(\bm{x}_4 = \times\)
--- all eight patterns on \(\bm{x}_1, \bm{x}_2, \bm{x}_4\), three inputs
shattered. The pairing supplies the third input for free. Hence

\begin{moral}
    \(\alpha \le B(3, 2)\).
\end{moral}

\subsection{Putting It All Together}

The two estimates add up, since \(2\alpha + \beta = (\alpha + \beta) +
\alpha\):
\begin{align*}
    B(4, 3) &= 2\alpha + \beta\\
    \alpha + \beta &\le B(3, 3)\\
    \alpha &\le B(3, 2)\\
    \implies B(4, 3) &\le B(3, 3) + B(3, 2) .
\end{align*}
Nothing in the argument used \(4\) or \(3\) specifically. For any \(N\) and
\(k\), take a largest table for \(B(N, k)\) and split it on the last input:
\begin{align*}
    B(N, k) &= 2\alpha + \beta\\
    \alpha + \beta &\le B(N - 1, k)\\
    \alpha &\le B(N - 1, k - 1)\\
    \implies B(N, k) &\le B(N - 1, k) + B(N - 1, k - 1) .
\end{align*}

\begin{lemma}[Inductive formula]\label{lem:B-recursion}
    For \(N \ge 2\) and \(k \ge 2\),
    \[
        B(N, k) \le B(N - 1, k) + B(N - 1, k - 1) .
    \]
\end{lemma}

\begin{proof}
    Let \(S\) be a set of \(B(N, k)\) vectors in \(\left\{ \circ, \times
    \right\}^N\) that shatters no \(k\) coordinates. Group its vectors by their
    first \(N - 1\) entries: \(\alpha\) prefixes appear with both values of the
    last entry and \(\beta\) with one, so \(\abs{S} = 2\alpha + \beta\). The
    \(\alpha + \beta\) distinct prefixes shatter no \(k\) of the first \(N - 1\)
    coordinates, or \(S\) would shatter them too; so \(\alpha + \beta \le
    B(N - 1, k)\). The \(\alpha\) paired prefixes shatter no \(k - 1\) of those
    coordinates, or \(S\) would shatter those \(k - 1\) together with the last
    one; so \(\alpha \le B(N - 1, k - 1)\). Add the two.
\end{proof}

Running the formula down the table from the known entries fills the region
below the diagonal with upper bounds (orange):

\begin{center}
    \solidrules
    \begin{tabular}{@{}cc|cccccc@{}}
        \toprule
        & & \multicolumn{6}{c}{\(k\)}\\
        & \(B(N, k)\) & \(1\) & \(2\) & \(3\) & \(4\) & \(5\) & \(6\)\\
        \midrule
        & \(1\) & \(1\) & \(2\) & \(2\) & \(2\) & \(2\) & \(2\)\\
        & \(2\) & \(1\) & \(3\) & \(4\) & \(4\) & \(4\) & \(4\)\\
        & \(3\) & \(1\) & \cellcolor{singlerowalt}\(4\)
        & \cellcolor{singlerowalt}\(7\) & \(8\) & \(8\) & \(8\)\\
        \(N\) & \(4\) & \(1\) & \textcolor{orange!85!black}{\(\le 5\)}
        & \cellcolor{singlerowalt}\(11\) & \(15\) & \(16\) & \(16\)\\
        & \(5\) & \(1\) & \textcolor{orange!85!black}{\(\le 6\)}
        & \textcolor{orange!85!black}{\(\le 16\)}
        & \textcolor{orange!85!black}{\(\le 26\)} & \(31\) & \(32\)\\
        & \(6\) & \(1\) & \textcolor{orange!85!black}{\(\le 7\)}
        & \textcolor{orange!85!black}{\(\le 22\)}
        & \textcolor{orange!85!black}{\(\le 42\)}
        & \textcolor{orange!85!black}{\(\le 57\)} & \(63\)\\
        \bottomrule
    \end{tabular}
\end{center}

The shaded cells show the pattern: each entry is at most the one directly
above plus the one above-left, as in Pascal's triangle. For instance,
\(B(4, 3) = 11 = 7 + 4\), so the formula is exact there, and
\(B(5, 3) \le 11 + 5 = 16\).

\begin{moral}
    We now have an upper bound of the bounding function.
\end{moral}

\subsection{Bounding Function: The Theorem}

The Pascal-like recursion suggests binomial coefficients, and indeed:

\begin{theorem}[Bounding function]\label{thm:bounding}
    For all \(N \ge 1\) and \(k \ge 1\),
    \[
        B(N, k) \le \underbrace{\sum_{i=0}^{k-1} \binom{N}{i}}_{\text{highest
        term } N^{k-1}} .
    \]
\end{theorem}

\begin{proof}
    Simple induction on \(N\), using the boundary cases and the inductive
    formula. Write \(\Phi(N, k)\) for the right-hand side.

    \emph{Boundary.} If \(k = 1\), \(B(N, 1) = 1 = \binom{N}{0}\) by
    \autoref{prop:B-basic}. If \(N = 1\) and \(k \ge 2\), \(B(1, k) = 2 =
    \binom{1}{0} + \binom{1}{1}\) (the sum stops at \(\binom{1}{1}\), all later
    terms being \(0\)).

    \emph{Induction.} For \(N \ge 2\), \(k \ge 2\), \autoref{lem:B-recursion}
    and the hypothesis for \(N - 1\) give
    \[
        B(N, k) \le \sum_{i=0}^{k-1} \binom{N-1}{i}
        + \sum_{i=0}^{k-2} \binom{N-1}{i}
        = \binom{N-1}{0} + \sum_{i=1}^{k-1}
        \left[ \binom{N-1}{i} + \binom{N-1}{i-1} \right]
        = \sum_{i=0}^{k-1} \binom{N}{i} ,
    \]
    the last step by Pascal's rule and \(\binom{N-1}{0} = \binom{N}{0}\).
\end{proof}

\begin{itemize}
    \item The proof is a simple induction using the boundary and the inductive
        formula.
    \item For fixed \(k\), \(B(N, k)\) is upper bounded by \(\mathrm{poly}(N)\),
        the leading term being \(\binom{N}{k-1} \approx N^{k-1} / (k-1)!\)
        \(\implies\) \(m_{\mathcal{H}}(N)\) is \(\mathrm{poly}(N)\) if a break
        point exists.
\end{itemize}

That settles \autoref{q:B-poly} and, through \(m_{\mathcal{H}}(N) \le B(N,
k)\), proves \autoref{conj:break}.

\begin{remark}
    `\(\le\)' can be `\(=\)' actually --- go play and prove it if you are a
    math lover! :-) (One way: exhibit a table of the right size, e.g.\ all
    vectors with fewer than \(k\) entries \(\times\). It has \(\sum_{i<k}
    \binom{N}{i}\) rows and cannot shatter \(k\) coordinates, since it never
    shows \(\times\) on all of them at once.) The inequality is also known as
    the Sauer--Shelah lemma.
\end{remark}

\subsection{The Three Break Points}

Returning to the hypothesis sets with a break point, \autoref{thm:bounding}
bounds each growth function from its break point alone (in orange, the bound):

\begin{center}
    \solidrules
    \begin{adjustbox}{max width=\textwidth}
        \begin{tabular}{@{}l|c|l|l@{}}
            \toprule
            \(\mathcal{H}\) & & \(m_{\mathcal{H}}(N)\) & break point\\
            \midrule
            positive rays & \dmark{o}\,\dmark{x}
            & \(N + 1 \;\textcolor{orange!85!black}{\le N + 1}\)
            & \(m_{\mathcal{H}}(2) = 3 < 2^2\): at \(2\)\\
            positive intervals & \dmark{o}\,\dmark{x}\,\dmark{o}
            & \(\tfrac{1}{2} N^2 + \tfrac{1}{2} N + 1 \;
            \textcolor{orange!85!black}{\le \tfrac{1}{2} N^2 + \tfrac{1}{2} N
            + 1}\)
            & \(m_{\mathcal{H}}(3) = 7 < 2^3\): at \(3\)\\
            2D perceptrons
            & \tikz[baseline=(current bounding box.center), x=0.3cm, y=0.3cm]{%
                \dmarkat{1}{2}{o}\dmarkat{0}{1}{x}\dmarkat{2}{1}{x}%
                \dmarkat{1}{0}{o}}
            & \(\textcolor{orange!85!black}{m_{\mathcal{H}}(N) = ? \le
            \tfrac{1}{6} N^3 + \tfrac{5}{6} N + 1}\)
            & \(m_{\mathcal{H}}(4) = 14 < 2^4\): at \(4\)\\
            \bottomrule
        \end{tabular}
    \end{adjustbox}
\end{center}

For rays and intervals the bound is exact: \(\binom{N}{0} + \binom{N}{1} = N +
1\), and \(\binom{N}{0} + \binom{N}{1} + \binom{N}{2} = \frac{1}{2} N^2 +
\frac{1}{2} N + 1\). For 2D perceptrons we still do not know
\(m_{\mathcal{H}}(N)\), but with \(k = 4\)
\[
    \sum_{i=0}^{3} \binom{N}{i}
    = 1 + N + \frac{N(N-1)}{2} + \frac{N(N-1)(N-2)}{6}
    = \frac{1}{6} N^3 + \frac{5}{6} N + 1 ,
\]
a polynomial, found without computing a single dichotomy beyond \(N = 4\).

\begin{moral}
    We can bound \(m_{\mathcal{H}}(N)\) by only one break point.
\end{moral}

\begin{exercise}[Fun Time]
    For 1D perceptrons (positive and negative rays), we know that
    \(m_{\mathcal{H}}(N) = 2N\). Let \(k\) be the minimum break point. Which of
    the following is \emph{not} true?
    \begin{enumerate}[(1)]
        \item \(k = 3\);
        \item for some integers \(N > 0\), \(m_{\mathcal{H}}(N) =
            \sum_{i=0}^{k-1} \binom{N}{i}\);
        \item for all integers \(N > 0\), \(m_{\mathcal{H}}(N) =
            \sum_{i=0}^{k-1} \binom{N}{i}\);
        \item for all integers \(N > 2\), \(m_{\mathcal{H}}(N) <
            \sum_{i=0}^{k-1} \binom{N}{i}\).
    \end{enumerate}
\end{exercise}

\begin{proof}[Reference Answer]
    (3). The proof is generally trivial by listing the definitions. (1) is the
    last Fun Time of \autoref{sec:break-point}. With \(k = 3\) the sum is
    \(1 + N + \frac{N(N-1)}{2}\), which equals \(2\), \(4\), \(7\) at
    \(N = 1, 2, 3\): so for (2), \(N = 1\) or \(2\) gives the equality, and
    \(N = 3\) (\(6 \ne 7\)) breaks (3). One thing to notice is (4): the
    difference \(1 + N + \frac{N(N-1)}{2} - 2N = \frac{(N-1)(N-2)}{2}\) is
    positive for \(N > 2\), so the upper bound can be `loose'.
\end{proof}

\section{A Pictorial Proof}\label{sec:pictorial}

\source{slides p.21--26}

The growth function is now known to be polynomial whenever there is a break
point. What remains is the `\(?\)' of \autoref{eq:wish}: may
\(m_{\mathcal{H}}\) take the place of \(M\) in the bound at all?

\subsection{BAD Bound for General \texorpdfstring{\(\mathcal{H}\)}{H}}

Because \(\mathcal{A}\) may return any hypothesis, the event to control is that
\emph{some} \(h\) is BAD. What we want is
\[
    \mathbb{P}\left[ \exists h \in \mathcal{H} \text{ s.t. }
    \abs{\Ein(h) - \Eout(h)} > \epsilon \right]
    \le 2 \cdot m_{\mathcal{H}}(N) \cdot \exp \left( -2 \epsilon^2 N \right) .
\]
What is actually true, when \(N\) is large enough, is
\[
    \mathbb{P}\left[ \exists h \in \mathcal{H} \text{ s.t. }
    \abs{\Ein(h) - \Eout(h)} > \epsilon \right]
    \le 2 \cdot 2 m_{\mathcal{H}}(2N) \cdot
    \exp \left( -2 \cdot \frac{1}{16} \epsilon^2 N \right) .
\]
The wish survives with worse constants: \(m_{\mathcal{H}}(2N)\) in place of
\(m_{\mathcal{H}}(N)\) (still polynomial), an extra factor \(2\), and
\(\epsilon^2 / 16\) in the exponent (still exponential decay). Next: a sketch
of the proof, in three steps, each of which removes one obstacle and costs one
of those constants.

\subsection{Step 1: Replace \texorpdfstring{\(\Eout\)}{Eout} by
\texorpdfstring{\(\Ein'\)}{Ein'}}

\[
    \frac{1}{2} \mathbb{P}\left[ \exists h \in \mathcal{H} \text{ s.t. }
    \abs{\Ein(h) - \Eout(h)} > \epsilon \right]
    \le \mathbb{P}\left[ \exists h \in \mathcal{H} \text{ s.t. }
    \abs{\Ein(h) - \textcolor{red!80!black}{\Ein'(h)}} > \frac{\epsilon}{2}
    \right]
\]

\begin{itemize}
    \item \(\Ein(h)\) takes finitely many values across \(\mathcal{H}\) (one
        per dichotomy), while \(\Eout(h)\) takes infinitely many --- so replace
        the evil \(\Eout\) first.
    \item How? Sample a \vocab{verification set} \(\mathcal{D}'\) of size \(N\)
        to calculate \(\Ein'\), the error on \(\mathcal{D}'\).
    \item BAD \(h\) of \(\Ein - \Eout\) \(\overset{\text{probably}}{\implies}\)
        BAD \(h\) of \(\Ein - \Ein'\).
\end{itemize}

\(\mathcal{D}'\) is never actually collected; it exists only inside the proof.

\begin{center}
    \begin{tikzpicture}[x=1.15cm, y=1.75cm]
        \fill[green!12] (0,0) -- plot[domain=0:2, samples=60]
            (\x, {exp(-(\x-2)^2/0.5)}) -- (2,0) -- cycle;
        \fill[red!15] (3.45,0) -- plot[domain=3.45:4.4, samples=30]
            (\x, {exp(-(\x-2)^2/0.5)}) -- (4.4,0) -- cycle;
        \draw[black!70] (0,0) -- (4.4,0);
        \draw[black!70] (0,0) -- (0,1.12);
        \draw[thick] plot[domain=0:4.4, samples=120]
            (\x, {exp(-(\x-2)^2/0.5)});
        \draw[black!70] (2,0) -- (2,1.0);
        \node[below, font=\small] at (2,0) {\(\Eout\)};
        \draw[-{Stealth[length=2mm]}, thick] (2,0.07) -- (1.15,0.07);
        \node[font=\small, text=green!50!black] at (1.45,0.19) {\(\Ein'\)};
        \node[font=\small, text=red!80!black] at (4.12,0.19) {\(\Ein\)};
        \node[rotate=90, font=\scriptsize, align=center] at (-0.32,0.56)
            {probability distribution\\of \(\Ein, \Ein'\)};
    \end{tikzpicture}
\end{center}

The picture is for one fixed \(h\). Both \(\Ein\) and \(\Ein'\) are averages of
\(N\) independent errors, so they share the same bell-shaped distribution
centred on \(\Eout\). Suppose \(\mathcal{D}\) is BAD, with \(\Ein\) out in the
red tail, more than \(\epsilon\) above \(\Eout\). The ghost sample is drawn
independently, so \(\Ein'\) lands in the green half, at or below \(\Eout\),
with probability about \(\frac{1}{2}\); and then \(\Ein\) and \(\Ein'\) are
more than \(\epsilon\) apart. So a BAD \(\mathcal{D}\) produces a large gap
\(\abs{\Ein - \Ein'}\) at least half the time, which is the \(\frac{1}{2}\) on
the left.

\begin{remark}
    The rigorous version uses a margin instead of symmetry. On the event that
    \(\abs{\Ein(h) - \Eout(h)} > \epsilon\) for some \(h\), fix one such \(h\)
    (it depends only on \(\mathcal{D}\)). By Hoeffding on \(\mathcal{D}'\),
    \(\abs{\Ein'(h) - \Eout(h)} \le \frac{\epsilon}{2}\) with probability at
    least \(1 - 2\exp(-\frac{1}{2} \epsilon^2 N)\), which is at least
    \(\frac{1}{2}\) once \(N\) is large enough --- the `\(N\) large enough' in
    the statement. When it happens, the triangle inequality gives
    \(\abs{\Ein(h) - \Ein'(h)} > \epsilon - \frac{\epsilon}{2} =
    \frac{\epsilon}{2}\). The halved \(\epsilon\) is the price of this margin.
\end{remark}

\begin{moral}
    The evil \(\Eout\) is removed by verification with `ghost data'.
\end{moral}

\subsection{Step 2: Decompose \texorpdfstring{\(\mathcal{H}\)}{H} by Kind}

\begin{align*}
    \text{BAD}
    &\le 2 \mathbb{P}\left[ \exists h \in \mathcal{H} \text{ s.t. }
    \abs{\Ein(h) - \Ein'(h)} > \frac{\epsilon}{2} \right]\\
    &\le 2 \textcolor{blue!75!black}{m_{\mathcal{H}}(2N)}
    \mathbb{P}\left[ \textcolor{blue!75!black}{\text{fixed } h} \text{ s.t. }
    \abs{\Ein(h) - \Ein'(h)} > \frac{\epsilon}{2} \right]
\end{align*}

\begin{itemize}
    \item \(\Ein\) is computed with \(\mathcal{D}\) and \(\Ein'\) with
        \(\mathcal{D}'\) --- and now \(m_{\mathcal{H}}\) comes to play.
    \item How? Both errors depend on \(h\) only through its values on the
        \(2N\) points of \(\mathcal{D} \cup \mathcal{D}'\), so the infinite
        \(\mathcal{H}\) becomes \(\abs{\mathcal{H}(\bm{x}_1, \dots , \bm{x}_N,
        \bm{x}_1', \dots , \bm{x}_N')}\) kinds, at most
        \(m_{\mathcal{H}}(2N)\).
    \item Union bound on the \(m_{\mathcal{H}}(2N)\) kinds.
\end{itemize}

The first line is Step 1 with both sides doubled; `BAD' abbreviates the
left-hand probability we started from. The price of this step is
\(m_{\mathcal{H}}(2N)\) rather than \(m_{\mathcal{H}}(N)\): the dichotomies
must be counted on the ghost points too.

\begin{center}
    \begin{adjustbox}{max width=\textwidth}
        \begin{tikzpicture}[x=1.4cm, y=1.4cm]
            % (a) one hypothesis: one small BAD region
            \draw[plotbox] (0,0) rectangle (1.8,2.8);
            \badblob[20]{0.90}{2.10}{red}
            \node[font=\scriptsize, align=center] at (0.90,1.35)
                {space of\\data sets};
            \node[circle, fill=black, inner sep=0.9pt] (d) at (0.55,0.50) {};
            \draw[-{Stealth[length=1.6mm]}] (-0.20,0.08) -- (d);
            \node[font=\scriptsize] at (-0.28,0.02) {\(\mathcal{D}\)};
            \node[font=\scriptsize, align=center, below] at (0.90,-0.08) {(a) Hoeffding\\inequality};
            % (b) union bound: every BAD region counted as if disjoint
            \begin{scope}[shift={(2.3,0)}]
                \draw[plotbox] (0,0) rectangle (1.8,2.8);
                \badblob[333]{0.31}{0.27}{red}
                \badblob[48]{0.84}{0.33}{cyan!80!black}
                \badblob[259]{1.45}{0.26}{blue!80!black}
                \badblob[214]{0.29}{0.63}{ForestGreen}
                \badblob[217]{0.84}{0.63}{Brown}
                \badblob[114]{1.41}{0.68}{pink}
                \badblob[31]{0.35}{1.05}{magenta}
                \badblob[113]{0.91}{1.03}{Goldenrod}
                \badblob[148]{1.41}{1.08}{violet}
                \badblob[292]{0.32}{1.41}{orange}
                \badblob[92]{0.87}{1.44}{teal}
                \badblob[96]{1.41}{1.42}{Purple}
                \badblob[32]{0.31}{1.78}{LimeGreen}
                \badblob[254]{0.91}{1.79}{RoyalBlue}
                \badblob[160]{1.50}{1.77}{red}
                \badblob[185]{0.33}{2.19}{cyan!80!black}
                \badblob[357]{0.87}{2.18}{blue!80!black}
                \badblob[153]{1.51}{2.11}{ForestGreen}
                \badblob[229]{0.33}{2.56}{Brown}
                \badblob[60]{0.87}{2.57}{pink}
                \badblob[175]{1.47}{2.49}{magenta}
                \node[font=\scriptsize, align=center, below] at (0.90,-0.08) {(b) union\\bound};
            \end{scope}
            % (c) now: the BAD regions of one kind overlap
            \begin{scope}[shift={(4.6,0)}]
                \draw[plotbox] (0,0) rectangle (1.8,2.8);
                \badblob[20]{0.99}{2.14}{red}
                \badblob[285]{0.95}{2.04}{cyan!80!black}
                \badblob[160]{0.64}{1.96}{blue!80!black}
                \badblob[254]{0.83}{2.13}{ForestGreen}
                \badblob[47]{0.76}{2.00}{Brown}
                \badblob[340]{1.05}{2.01}{pink}
                \badblob[158]{1.11}{2.11}{magenta}
                \badblob[228]{0.63}{1.80}{Goldenrod}
                \badblob[342]{0.87}{2.15}{violet}
                \badblob[181]{0.71}{2.24}{orange}
                \badblob[30]{0.93}{2.09}{teal}
                \badblob[126]{0.92}{2.13}{Purple}
                \badblob[254]{0.65}{2.18}{LimeGreen}
                \badblob[281]{1.03}{2.10}{RoyalBlue}
                \node[font=\scriptsize, align=center, below] at (0.90,-0.08) {(c) now\\\strut};
            \end{scope}
        \end{tikzpicture}
    \end{adjustbox}
\end{center}

Each box is the space of all possible data sets, and each coloured scribble the
BAD region of one hypothesis --- the data sets on which it is fooled. (a)~Hoeffding
says one hypothesis's region is small, so a random \(\mathcal{D}\) usually
misses it. (b)~The union bound adds up all the regions as if they were
disjoint, so with many hypotheses they seem to fill the whole space. (c)~In
truth, hypotheses of the same kind are BAD on nearly the same data sets, and
their regions pile on top of one another; counting kinds instead of hypotheses
counts each pile once.

\begin{moral}
    Use \(m_{\mathcal{H}}(2N)\) to calculate the BAD-overlap properly.
\end{moral}

\subsection{Step 3: Use Hoeffding without Replacement}

\begin{align*}
    \text{BAD}
    &\le 2 m_{\mathcal{H}}(2N) \mathbb{P}\left[ \text{fixed } h \text{ s.t. }
    \abs{\Ein(h) - \Ein'(h)} > \frac{\epsilon}{2} \right]\\
    &\le 2 m_{\mathcal{H}}(2N) \cdot 2 \exp \left( -2
    \left( \frac{\epsilon}{4} \right)^2 N \right)
\end{align*}

\begin{itemize}
    \item Consider a bin of the \(2N\) examples; choose \(N\) of them for
        \(\Ein\), and leave the others for \(\Ein'\). Since the average of the
        two halves is the bin's own error rate,
        \[
            \abs{\Ein - \Ein'} > \frac{\epsilon}{2}
            \iff
            \abs{\Ein - \frac{\Ein + \Ein'}{2}} > \frac{\epsilon}{4} .
        \]
    \item So? It is just a `smaller bin', a `smaller \(\epsilon\)', and
        Hoeffding without replacement.
\end{itemize}

\begin{center}
    \begin{tikzpicture}
        \begin{scope}[scale=1.9]
            \foreach \i in {0,...,9}{
                \pgfmathsetmacro{\xx}{-0.324+0.072*\i}
                \ifodd\i
                    \node[morange=3.8pt] at (\xx,0.050) {};
                    \node[mgreen=3.8pt] at (\xx+0.036,0.118) {};
                \else
                    \node[mgreen=3.8pt] at (\xx,0.050) {};
                    \node[morange=3.8pt] at (\xx+0.036,0.118) {};
                \fi}
        \end{scope}
        \draw[binwall, scale=1.9] \binpath;
        \node[font=\small\bfseries, below] at (0,-0.10) {small bin};
        \draw[scoop] (0.05,1.95) .. controls (0.85,2.75) and (2.20,2.45) ..
            (2.55,1.10);
        \node[font=\small\bfseries] at (1.40,2.70) {sample for \(\Ein\)};
        \foreach \i/\c in {0/mgreen,1/mgreen,2/morange,3/mgreen,4/morange,
            5/mgreen,6/mgreen,7/morange,8/morange,9/mgreen}{
            \node[\c=4.5pt] at (2.05+0.17*\i,0.85) {};}
    \end{tikzpicture}
\end{center}

Everything is now conditioned on the \(2N\) points in the bin, and \(h\) is
fixed. Which \(N\) of them are \(\mathcal{D}\) and which are \(\mathcal{D}'\)
is a uniformly random split, so \(\Ein\) is the orange fraction of a sample of
\(N\) drawn \emph{without} replacement from a bin whose orange fraction is
\(\frac{\Ein + \Ein'}{2}\). Hoeffding's inequality holds for sampling without
replacement as well (Hoeffding proved it in that form), with the tolerance now
\(\frac{\epsilon}{4}\): that gives \(2\exp(-2(\epsilon/4)^2 N)\).

\begin{moral}
    Use Hoeffding after zooming to a fixed \(h\).
\end{moral}

\subsection{That's All!}

Collecting the constants, \(2 m_{\mathcal{H}}(2N) \cdot 2 \exp(-2
\frac{\epsilon^2}{16} N) = 4 m_{\mathcal{H}}(2N) \exp(-\frac{1}{8} \epsilon^2
N)\).

\begin{theorem}[Vapnik--Chervonenkis (VC) bound]\label{thm:vc}
    For \(N\) large enough,
    \[
        \mathbb{P}\left[ \exists h \in \mathcal{H} \text{ s.t. }
        \abs{\Ein(h) - \Eout(h)} > \epsilon \right]
        \le 4 m_{\mathcal{H}}(2N) \exp \left( -\frac{1}{8} \epsilon^2 N
        \right) .
    \]
\end{theorem}

\begin{proof}[Proof sketch]
    \leavevmode
    \begin{itemize}
        \item replace \(\Eout\) by \(\Ein'\) (Step 1);
        \item decompose \(\mathcal{H}\) by kind (Step 2);
        \item use Hoeffding without replacement (Step 3). \qedhere
    \end{itemize}
\end{proof}

Since the event covers every \(h\), it covers the \(g\) returned by any
\(\mathcal{A}\), exactly as \autoref{thm:finite-hoeffding} did for finite
\(\mathcal{H}\). Applied to 2D perceptrons:

\begin{itemize}
    \item break point? \(4\);
    \item \(m_{\mathcal{H}}(N)\)? \(O(N^3)\), by \autoref{thm:bounding}, so
        \(m_{\mathcal{H}}(2N)\) is \(O(N^3)\) as well, and the bound goes to
        \(0\) as \(N\) grows.
\end{itemize}

\begin{moral}
    Learning with 2D perceptrons is feasible! :-)
\end{moral}

\begin{exercise}[Fun Time]
    For positive rays, \(m_{\mathcal{H}}(N) = N + 1\). Plug it into the VC
    bound for \(\epsilon = 0.1\) and \(N = 10000\). What is the VC bound of
    BAD events?
    \[
        \mathbb{P}\left[ \exists h \in \mathcal{H} \text{ s.t. }
        \abs{\Ein(h) - \Eout(h)} > \epsilon \right]
        \le 4 m_{\mathcal{H}}(2N) \exp \left( -\frac{1}{8} \epsilon^2 N
        \right)
    \]
    \begin{enumerate*}[(1)]
        \item \(2.77 \times 10^{-87}\); \item \(5.54 \times 10^{-83}\);
        \item \(2.98 \times 10^{-1}\); \item \(2.29 \times 10^{2}\).
    \end{enumerate*}
\end{exercise}

\begin{proof}[Reference Answer]
    (3). Simple calculation:
    \[
        4 \cdot (2 \cdot 10000 + 1) \cdot \exp \left( -\frac{1}{8} (0.1)^2
        \cdot 10000 \right)
        = 80004 \, e^{-12.5} \approx 0.298 .
    \]
    Note that the BAD probability bound is not very small even with \(10000\)
    examples --- the constants lost in the three steps are expensive.
\end{proof}

\hrulebar

\subsection*{Summary}

\begin{description}
    \item[Restriction of Break Point] A break point `breaks' consequent points.
    \item[Bounding Function: Basic Cases] \(B(N, k)\) bounds
        \(m_{\mathcal{H}}(N)\) with break point \(k\).
    \item[Bounding Function: Inductive Cases] \(B(N, k)\) is
        \(\mathrm{poly}(N)\).
    \item[A Pictorial Proof] \(m_{\mathcal{H}}(N)\) can replace \(M\) with a
        few changes.
\end{description}

The chapter closed both gaps left by the last one. A break point \(k\) caps
\(m_{\mathcal{H}}(N)\) by \(\sum_{i<k} \binom{N}{i} = O(N^{k-1})\), and the VC
bound puts \(m_{\mathcal{H}}(2N)\) where \(M\) used to be; so any
\(\mathcal{H}\) with a break point generalises, however infinite it is. Next:
how to `use' the break point --- the minimum break point turns out to be the
one number that summarises an \(\mathcal{H}\)'s capacity.
